\begin{lemma}\label{lemma:R_4n_add_1}\lean{R_4n_add_1}\leanok
For all $n > 1$, $R(4n + 1) = R(n)$.
\end{lemma}
\begin{proof}\leanok
We compute $3(4n+1) + 1 = 12n + 4 = 4(3n+1) = 2^2 \cdot (3n+1)$. Since multiplying by $2^2$ does not change the odd part, we have $R(4n+1) = \text{odd part of } 2^2(3n+1) = \text{odd part of } (3n+1) = R(n)$.
\end{proof}
\begin{lemma}\label{lemma:R_iter_4n_add_1}\lean{R_iter_4n_add_1}\leanok
For all $n > 1$ and $i > 0$, $R^i(4n + 1) = R^i(n)$.
\end{lemma}
\begin{proof}\leanok
For $i = k + 1$, we have
$R^{k+1}(4n+1) = R^k(R(4n+1)) = R^k(R(n)) = R^{k+1}(n)$,
where the middle equality follows from Lemma~\ref{lemma:R_4n_add_1}.
\end{proof}

\begin{lemma}\label{lemma:rcs_ge}\lean{rcs_ge}\leanok
For any odd $n \in \mathbb{N}$, $2 R(n) \le 3n + 1$.
\end{lemma}
\begin{proof}\leanok
By definition, $R(n)$ is the odd part of $3n + 1$, meaning $3n + 1 = 2^k \cdot R(n)$ for some $k \in \mathbb{N}$. Since $n$ is odd, $3n + 1$ is even, which implies $k \ge 1$. Thus $2^k \ge 2$, and we have $3n + 1 \ge 2 R(n)$.
\end{proof}

\begin{definition}\label{def:IsPrimitive4x1}\lean{IsPrimitive4x1}\leanok
An odd number $n > 1$ is called a \emph{primitive} at level $i > 1$ if
\begin{enumerate}
\item $R^i(n) = 1$,
\item there is no odd $k$ with $4k + 1 = n$ and $R^i(k) = 1$.
\end{enumerate}
\end{definition}

\begin{lemma}\label{lemma:is_primitive_iff_mod_8_ne_5}\lean{is_primitive_iff_mod_8_ne_5}\leanok
Let $n > 1$ be odd with $n \neq 5$, let $i > 1$, and suppose $R^i(n) = 1$.
Then $n$ is a primitive at level $i$ if and only if $n \not\equiv 5 \pmod{8}$.
\end{lemma}
\begin{proof}\leanok
($\Rightarrow$) Suppose $n$ is primitive and $n \equiv 5 \pmod{8}$.
Write $n = 8k + 5 = 4(2k+1) + 1$. Then $2k+1$ is odd and $2k + 1 > 1$ (since $n \neq 5$ implies $k \geq 1$).
By Lemma~\ref{lemma:R_iter_4n_add_1}, $R^i(2k+1) = R^i(4(2k+1)+1) = R^i(n) = 1$,
so $2k+1$ is an odd predecessor of $n$ via the $4k+1$ map with the same step count, contradicting primitivity.

($\Leftarrow$) Suppose $n \not\equiv 5 \pmod{8}$ and there exists an odd $k$ with $4k + 1 = n$ and $R^i(k) = 1$.
Since $k$ is odd, write $k = 2m + 1$; then $n = 4(2m+1) + 1 = 8m + 5 \equiv 5 \pmod{8}$, a contradiction.
\end{proof}

\begin{lemma}\label{lemma:primitive_ancestor}\lean{primitive_ancestor}\leanok
Let $x > 1$ be odd with $R(x) \neq 1$. Then there exists an odd $n > 1$ with $R(n) = R(x)$ and $n \not\equiv 5 \pmod{8}$.
\end{lemma}
\begin{proof}\leanok
We proceed by strong induction on $x$.
If $x \not\equiv 5 \pmod{8}$, take $n = x$.

If $x \equiv 5 \pmod{8}$, write $x = 8k + 5 = 4(2k+1) + 1$.
We must have $2k + 1 > 1$, since $k = 0$ gives $x = 5$ and $R(5) = 1$, contradicting $R(x) \neq 1$.
By Lemma~\ref{lemma:R_4n_add_1}, $R(x) = R(2k+1)$, so $R(2k+1) \neq 1$.
Since $2k + 1 < x$, the induction hypothesis yields an odd $n > 1$ with $R(n) = R(2k+1) = R(x)$ and $n \not\equiv 5 \pmod{8}$.
\end{proof}

\begin{lemma}\label{lemma:odd_node_generates_primitive}\lean{odd_node_generates_primitive}\leanok
Let $y > 1$ be odd with $y \not\equiv 0 \pmod{3}$ and $R^i(y) = 1$. Then there exists a primitive $n$ at level $i + 1$.
\end{lemma}
\begin{proof}\leanok
Since $y$ is odd, positive, and not divisible by~3, there exists an odd $x > 1$ with $R(x) = y$ (by the surjectivity of $R$ onto odd numbers not divisible by~3).
In particular $R(x) \neq 1$ since $y > 1$.

By Lemma~\ref{lemma:primitive_ancestor}, there exists an odd $n > 1$ with $R(n) = R(x) = y$ and $n \not\equiv 5 \pmod{8}$.
Then $R^{i+1}(n) = R^i(R(n)) = R^i(y) = 1$.
Since $n \neq 5$ (as $5 \equiv 5 \pmod{8}$), Lemma~\ref{lemma:is_primitive_iff_mod_8_ne_5} gives that $n$ is a primitive at level $i + 1$.
\end{proof}

\begin{definition}\label{def:reduce4x1}\lean{reduce4x1}\leanok
The function $\text{reduce4x1}: \mathbb{N} \to \mathbb{N}$ is defined recursively by
\[ \text{reduce4x1}(n) = \begin{cases} \text{reduce4x1}((n-1)/4) & \text{if } n \equiv 5 \pmod{8}, \\ n & \text{otherwise.} \end{cases} \]
\end{definition}

\begin{lemma}\label{lemma:rcs_reduce4x1}\lean{rcs_reduce4x1}\leanok
For any $n \in \mathbb{N}$, $R(\text{reduce4x1}(n)) = R(n)$.
\end{lemma}
\begin{proof}\leanok
We proceed by strong induction on $n$. If $n \not\equiv 5 \pmod{8}$, then $\text{reduce4x1}(n) = n$ and the result is trivial.
If $n \equiv 5 \pmod{8}$, then $n = 4k+1$ for some odd $k = (n-1)/4$.
By the definition of $\text{reduce4x1}$, $\text{reduce4x1}(n) = \text{reduce4x1}(k)$.
By induction, $R(\text{reduce4x1}(k)) = R(k)$.
By Lemma~\ref{lemma:R_4n_add_1}, $R(n) = R(4k+1) = R(k)$.
Thus $R(\text{reduce4x1}(n)) = R(n)$.
\end{proof}

\begin{lemma}\label{lemma:primitives_from_distinct_generators_ne}\lean{primitives_from_distinct_generators_ne}\leanok
Let $x_1, x_2 \in \mathbb{N}$ and $y_1, y_2$ be their respective images under $R$, with $y_1 \neq y_2$. Then $\text{reduce4x1}(x_1) \neq \text{reduce4x1}(x_2)$.
\end{lemma}
\begin{proof}\leanok
Suppose for contradiction that $\text{reduce4x1}(x_1) = \text{reduce4x1}(x_2)$.
Then by Lemma~\ref{lemma:rcs_reduce4x1},
$y_1 = R(x_1) = R(\text{reduce4x1}(x_1)) = R(\text{reduce4x1}(x_2)) = R(x_2) = y_2$,
contradicting $y_1 \neq y_2$.
\end{proof}

\begin{lemma}\label{lemma:exists_large_not_div3_from_seed}\lean{exists_large_not_div3_from_seed}\leanok
Given a level $m \ge 1$, a seed $y_0 > 1$ that is odd and satisfies $R^m(y_0) = 1$, and $B \in \mathbb{N}$, there exists an odd $y > B$ with $y > 1$, $y \not\equiv 0 \pmod{3}$, and $R^m(y) = 1$.
\end{lemma}
\begin{proof}\leanok
We first show by induction on $B$ that there exists an odd $y > B$ with $y > 1$ and $R^m(y) = 1$.
The base case $B=0$ is given by $y_0$.
For the inductive step, if $y > B$ satisfies the conditions, let $y' = 4y + 1$. Then $y' > 4B + 1 > B+1$, $y'$ is odd, $y' > 1$, and $R^m(y') = R^m(y) = 1$ by Lemma~\ref{lemma:R_iter_4n_add_1}.

Now, fixed such a $y > B$ with $R^m(y) = 1$ and $y$ odd, $y > 1$. If $y \not\equiv 0 \pmod{3}$ we are done.
Otherwise, if $y \equiv 0 \pmod{3}$, consider $y' = 4y + 1$.
Then $y' > y > B$, $y'$ is odd, $y' > 1$, and $R^m(y') = 1$.
Finally, $4y + 1 \equiv 4(0) + 1 \equiv 1 \pmod{3}$, so $y' \not\equiv 0 \pmod{3}$.
\end{proof}

\begin{lemma}\label{lemma:infinite_primitives}\lean{infinite_primitives}\leanok
For every level $m \ge 2$, there are infinitely many primitive numbers at level $m$.
\end{lemma}
\begin{proof}\leanok
We proceed by induction on $m \ge 2$.
Base case $m=2$: Let $B \in \mathbb{N}$. We apply Lemma~\ref{lemma:exists_large_not_div3_from_seed} starting from the seed $y_0 = 5$ at level 1 (with $R(5)=1$) to find an odd $y > 2B + 2$ with $y \not\equiv 0 \pmod{3}$ and $R(y) = 1$. By Lemma~\ref{lemma:odd_node_generates_primitive}, there exists a primitive $n$ at level 2 such that $R(n) = y$. Since $2R(n) \le 3n + 1$ (by Lemma~\ref{lemma:rcs_ge}), $n \ge (2y - 1)/3 > \frac{4B+3}{3} > B$. Thus $n > B$ is the required primitive.

Inductive step $m \to m+1$: Suppose there is a primitive $p$ at level $m$. By Lemma~\ref{lemma:exists_large_not_div3_from_seed}, there exists an odd $y > 2B + 2$ with $y > 1$, $y \not\equiv 0 \pmod{3}$, and $R^m(y) = 1$. By Lemma~\ref{lemma:odd_node_generates_primitive}, there is a primitive $n$ at level $m+1$ with $R(n) = y$. As in the base case, the bound $2R(n) \le 3n + 1$ ensures $n > B$.
\end{proof}
